% Options for packages loaded elsewhere
\PassOptionsToPackage{unicode}{hyperref}
\PassOptionsToPackage{hyphens}{url}
\PassOptionsToPackage{dvipsnames,svgnames*,x11names*}{xcolor}
%
\documentclass[
  10pt,
]{article}
\usepackage{amsmath,amssymb}
\usepackage{lmodern}
\usepackage{iftex}
\ifPDFTeX
  \usepackage[T1]{fontenc}
  \usepackage[utf8]{inputenc}
  \usepackage{textcomp} % provide euro and other symbols
\else % if luatex or xetex
  \usepackage{unicode-math}
  \defaultfontfeatures{Scale=MatchLowercase}
  \defaultfontfeatures[\rmfamily]{Ligatures=TeX,Scale=1}
  \setmainfont[]{Charter}
  \setmonofont[Scale=0.80]{Menlo}
\fi
% Use upquote if available, for straight quotes in verbatim environments
\IfFileExists{upquote.sty}{\usepackage{upquote}}{}
\IfFileExists{microtype.sty}{% use microtype if available
  \usepackage[]{microtype}
  \UseMicrotypeSet[protrusion]{basicmath} % disable protrusion for tt fonts
}{}
\makeatletter
\@ifundefined{KOMAClassName}{% if non-KOMA class
  \IfFileExists{parskip.sty}{%
    \usepackage{parskip}
  }{% else
    \setlength{\parindent}{0pt}
    \setlength{\parskip}{6pt plus 2pt minus 1pt}}
}{% if KOMA class
  \KOMAoptions{parskip=half}}
\makeatother
\usepackage{xcolor}
\IfFileExists{xurl.sty}{\usepackage{xurl}}{} % add URL line breaks if available
\IfFileExists{bookmark.sty}{\usepackage{bookmark}}{\usepackage{hyperref}}
\hypersetup{
  colorlinks=true,
  linkcolor={RoyalBlue},
  filecolor={Maroon},
  citecolor={Blue},
  urlcolor={RoyalBlue},
  pdfcreator={LaTeX via pandoc}}
\urlstyle{same} % disable monospaced font for URLs
\usepackage[margin=0.9in]{geometry}
\usepackage{longtable,booktabs,array}
\usepackage{calc} % for calculating minipage widths
% Correct order of tables after \paragraph or \subparagraph
\usepackage{etoolbox}
\makeatletter
\patchcmd\longtable{\par}{\if@noskipsec\mbox{}\fi\par}{}{}
\makeatother
% Allow footnotes in longtable head/foot
\IfFileExists{footnotehyper.sty}{\usepackage{footnotehyper}}{\usepackage{footnote}}
\makesavenoteenv{longtable}
\setlength{\emergencystretch}{3em} % prevent overfull lines
\providecommand{\tightlist}{%
  \setlength{\itemsep}{0pt}\setlength{\parskip}{0pt}}
\setcounter{secnumdepth}{-\maxdimen} % remove section numbering
\usepackage{newunicodechar}
\newunicodechar{→}{\ensuremath{\rightarrow}}
\newunicodechar{←}{\ensuremath{\leftarrow}}
\newunicodechar{↔}{\ensuremath{\leftrightarrow}}
\newunicodechar{⇒}{\ensuremath{\Rightarrow}}
\newunicodechar{≈}{\ensuremath{\approx}}
\newunicodechar{≤}{\ensuremath{\leq}}
\newunicodechar{≥}{\ensuremath{\geq}}
\newunicodechar{≠}{\ensuremath{\neq}}
\newunicodechar{×}{\ensuremath{\times}}
\newunicodechar{±}{\ensuremath{\pm}}
\newunicodechar{∈}{\ensuremath{\in}}
\newunicodechar{≪}{\ensuremath{\ll}}
\ifLuaTeX
  \usepackage{selnolig}  % disable illegal ligatures
\fi

\title{SDR Platform — Applications Introduction}
\author{}
\date{}

\begin{document}
\maketitle

{
\hypersetup{linkcolor=}
\setcounter{tocdepth}{2}
\tableofcontents
}
\hypertarget{sdr-platform--applications-introduction}{%
\section{SDR Platform --- Applications
Introduction}\label{sdr-platform--applications-introduction}}

A single software-defined-radio physical layer (\texttt{sdr\_system},
C++/UHD) carries several very different \textbf{applications}. This
document introduces each one in detail: what it computes, how it uses
the radio, the exact signal chain it exercises, the hardware topology it
runs on, and its current status. It is written for a collaborator coming
to the code fresh (OSU experiments / UnionLabs deep-dive).

Two applications exist today, and they sit at \textbf{opposite ends of
the design space}:

\begin{longtable}[]{@{}lll@{}}
\toprule
& \textbf{Application 1 --- Reliable Digital Link} & \textbf{Application
2 --- Over-the-Air Computation} \\
\midrule
\endhead
Paradigm & decode-and-forward (one packet, one receiver) & analog
superposition (many senders, one function) \\
Transmitters active at once & \textbf{1} (per slot) & \textbf{N}
(simultaneously) \\
What the receiver produces & the \textbf{delivered bytes} + an ACK & an
\textbf{aggregate} (the sum) of all senders \\
Receiver work & full sync + demod + FEC decode + CRC & CSI + a fixed
linear combine \\
Hosts & MARL random access, OSU, Federated Learning & STLC digital
AirComp (Lee--Lee--Jung, WCL 2026) \\
Status & validated over the air & design mapped; new blocks to build \\
\bottomrule
\end{longtable}

The rest of the document is: the shared substrate both applications
stand on (§1), then each application in full (§2, §3), then how they
unify under one abstraction (§4), and a reference appendix (§5).

\begin{center}\rule{0.5\linewidth}{0.5pt}\end{center}

\hypertarget{1-the-shared-substrate}{%
\subsection{1. The shared substrate}\label{1-the-shared-substrate}}

Everything below the application logic is common. An application never
re-implements DSP; it either drives the monolithic engine through a CLI
wrapper, or composes the PHY as GNU-Radio-style blocks in Python.

\hypertarget{11-the-phy-engine-sdr_system-c}{%
\subsubsection{\texorpdfstring{1.1 The PHY engine (\texttt{sdr\_system},
C++)}{1.1 The PHY engine (sdr\_system, C++)}}\label{11-the-phy-engine-sdr_system-c}}

A complete two-radio digital link. Full detail lives in
\texttt{../SYSTEM\_REFERENCE.pdf}; the parts the applications rely on:

\begin{verbatim}
TX:  message -> chunk -> CRC-16 -> [FEC encode] -> bits->symbols (modulator)
        |-- single-carrier: RRC pulse-shape -> preamble -> USRP TX
        |-- OFDM:            QAM -> IFFT+CP frame -> USRP TX
RX:  USRP RX -> energy gate -> AGC -> sync
        |-- single-carrier: matched filter -> timing -> [equalize] -> demod
        |-- OFDM:            Schmidl-Cox + CFO -> FFT -> 1-tap eq -> pilot CPE -> demod
     -> [FEC decode] -> CRC-16 -> (ARQ: ACK if OK) -> reassemble
\end{verbatim}

\begin{itemize}
\tightlist
\item
  \textbf{Modulation:} BPSK, QPSK, 8-PSK, 16/32/64/128/256-QAM, and
  differential DBPSK / DQPSK / 8-DPSK / PI4-QPSK.
\item
  \textbf{Waveform:} single-carrier (RRC, roll-off 0.25) or OFDM
  (per-subcarrier pilots that track a free-running carrier-frequency
  offset).
\item
  \textbf{FEC:} rate-1/2 convolutional (K=7) + Viterbi, \textbf{LDPC}
  (IRA/staircase, normalized min-sum belief propagation), and
  \textbf{turbo} (punctured PCCC, iterative max-log-MAP BCJR). Hard or
  soft decision; tunable block size \texttt{k}, iterations, and
  normalization.
\item
  \textbf{Reliability:} CRC-16, stop-and-wait ARQ, ACK carried over TCP
  or over RF.
\item
  \textbf{Channel awareness:} energy sensing, \texttt{freq\_scan.py},
  live EVM and pre-/post-FEC BER.
\end{itemize}

\textbf{Free-running-LO fact that shapes both applications:} the two
radios' local oscillators run independently. Coherent single-carrier
QPSK then fails (\textasciitilde50 \% BER) because the carrier phase
spins; the two remedies are \textbf{differential} modulation
(DQPSK/DBPSK --- no absolute phase reference needed) or \textbf{OFDM}
(pilots track the offset). This is why differential modes are
first-class, and it directly informs the AirComp design in §3.

\hypertarget{12-the-block-api-pyphy}{%
\subsubsection{\texorpdfstring{1.2 The block API
(\texttt{pyphy})}{1.2 The block API (pyphy)}}\label{12-the-block-api-pyphy}}

Every DSP stage is also exposed as a numpy-in / numpy-out function via
pybind11, so an application can compose the chain in Python and insert
its own operations between stages (exactly the GNU Radio model):

\begin{verbatim}
framing : frame / unframe
fec     : fec_encode / fec_decode / fec_decode_soft        (conv|ldpc|turbo)
mod     : modulate / demodulate / soft_llr
sc      : rrc_tx / rrc_rx
sync    : preamble / acq / cfo_correct / phase_correct
ofdm    : ofdm_mod / ofdm_demod
radio   : Radio('tx'|'rx', ...).transmit(wave) / .capture(n)      (WITH_UHD build)
\end{verbatim}

The same C++ the binary uses, so behaviour is identical. This block API
is the substrate Application 2 will be built on (custom transmit math +
raw multi-antenna capture + custom receiver estimation are impossible in
the monolithic binary but natural here).

\hypertarget{13-the-abstraction-layer-in-progress}{%
\subsubsection{1.3 The abstraction layer (in
progress)}\label{13-the-abstraction-layer-in-progress}}

For the UnionLabs collaboration we are standardizing how an algorithm
plugs into the PHY: an \texttt{SdrApp} implements three methods ---
\texttt{next\_payload()\ -\textgreater{}\ float32{[}{]}} (PHY grabs it
to transmit), \texttt{on\_payload(float32{[}{]})} (received data handed
back), \texttt{on\_result(ack)} (delivered / collision, i.e. the reward)
--- and declares one \texttt{PayloadSpec(dtype,\ shape)}. Adding an
application becomes "3 methods + 1 spec"; nothing below the seam
changes. §4 revisits this once both applications are on the table,
because Application 2 stretches the contract.

\begin{center}\rule{0.5\linewidth}{0.5pt}\end{center}

\hypertarget{2-application-1--reliable-digital-link}{%
\subsection{2. Application 1 --- Reliable Digital
Link}\label{2-application-1--reliable-digital-link}}

\hypertarget{21-what-it-is}{%
\subsubsection{2.1 What it is}\label{21-what-it-is}}

A classic \textbf{decode-and-forward} link: one device transmits a
packet, one access point recovers the exact bytes, and an ACK closes the
loop. The "application" living on top decides \emph{what} to send and
\emph{what the ACK means}. Three algorithms share this link, in two
archetypes:

\begin{itemize}
\tightlist
\item
  \textbf{Data-transfer} --- the payload \emph{is} the information (a
  float32 vector). Uses \texttt{next\_payload()} +
  \texttt{on\_payload()}. Example: \textbf{Federated Learning}.
\item
  \textbf{Control / random-access} --- the payload is fixed; the
  \emph{outcome} is the signal. Uses \texttt{next\_payload()} +
  \texttt{on\_result(ack)}, where the \textbf{ACK boolean is the
  reinforcement-learning reward}. Examples: \textbf{MARL} random access
  and \textbf{OSU}.
\end{itemize}

The full PHY chain of §1.1 is exercised end-to-end: framing, FEC,
modulation, sync, demodulation, decoding, CRC, ARQ.

\hypertarget{22-marl-random-access-control-archetype}{%
\subsubsection{2.2 MARL random access (control
archetype)}\label{22-marl-random-access-control-archetype}}

Multi-agent reinforcement learning over a contended channel. Topology is
\textbf{N transmitting agents -\textgreater{} 1 receiving access point}.

\begin{itemize}
\tightlist
\item
  \textbf{The reward is the ACK.} Each agent fires a single-shot burst
  (\texttt{-\/-max-attempts\ 1}); if its frame decodes and the AP routes
  back an ACK, that is a success (reward), otherwise a timeout means
  collision/loss. The policy --- not the PHY --- owns retransmission.
\item
  \textbf{Per-agent ACK routing (1 RX -\textgreater{} N TX).} You do not
  need N receivers. The agent id is carried in payload byte 0; the AP
  runs the C++ decoder with \texttt{role\ rx\ -\/-marl-report}, which
  prints one
  \texttt{{[}BURST{]}\ id=\ldots{}\ idx\ tot\ nbytes\ hex=\ldots{}} line
  per CRC-OK burst. \texttt{ap\_multi.py} parses those lines and routes
  an ACK to the agent whose frame decoded. A collision produces no line,
  hence no ACK, hence a natural timeout.
\item
  \textbf{Slotted time (\texttt{slot\_sync.py}).} So contention is well
  defined, a standalone TCP slot clock broadcasts
  \texttt{SLOT\ \textless{}k\textgreater{}} every \texttt{-\/-slot-ms};
  agents block on \texttt{wait\_slot()} at the top of each decision. The
  AP arbitrates at each slot boundary: \textgreater=2 intents
  -\textgreater{} COLLISION (no ACK to anyone, even if one happened to
  decode via capture); exactly 1 intent that decoded -\textgreater{}
  real ACK; 1 intent that did not decode -\textgreater{} real loss; 0
  -\textgreater{} idle. ACKs are emitted only at arbitration, so
  collisions are deterministic rather than dependent on RF-overlap
  timing.
\item
  \textbf{Warm radios.} Both source and sink stay warm (LOs run
  continuously) so the CFO is stable and trackable --- this lets
  coherent QPSK win (85 \% vs DQPSK 62 \% delivery on the directional
  link) instead of being forced onto differential by cold-LO jumps.
\item
  \textbf{Observation vs payload.} The agent's RL state
  \texttt{{[}AoI,\ queue,\ channel-busy{]}} (float32) is \emph{local}
  and is \textbf{not} transported --- only the fixed burst crosses the
  air. This is the clean separation the abstraction layer formalizes.
\end{itemize}

\hypertarget{23-osu-control-archetype}{%
\subsubsection{2.3 OSU (control
archetype)}\label{23-osu-control-archetype}}

OSU is the same control loop as MARL with a different policy swapped in.
It observes its state, decides transmit-or-defer, fires one burst
through the adapter, and reads back delivered/collision as reward. It is
the reference application we are porting onto the \texttt{SdrApp}
contract first; the goal of the deep-dive is to make "bring your own
policy" a 3-methods-and-a-spec exercise.

\hypertarget{24-federated-learning-data-transfer-archetype}{%
\subsubsection{2.4 Federated Learning (data-transfer
archetype)}\label{24-federated-learning-data-transfer-archetype}}

FedAvg of an MNIST classifier over the radio (\texttt{fl.py} /
\texttt{fl\_core.py}, numpy-only, torch-free).

\begin{itemize}
\tightlist
\item
  \textbf{What crosses the air:} model updates as float32 vectors. In
  the default \textbf{compressed} mode all nodes share an init seed and
  exchange only \textbf{top-k sparse deltas} with error feedback both
  directions; the server FedAvg-averages the sparse client deltas and
  broadcasts the sparse aggregate. This is \textasciitilde20 KB/delta
  versus \textasciitilde200 KB for the full model. Mock-validated live
  on MNIST: \textbf{0.71 -\textgreater{} 0.93 accuracy in 20 rounds, 2
  clients} (beating full-model 0.925).
\item
  \textbf{Per-direction channel.} Because the rig has an RX-only N210
  and TX-only B210s, the only RF path is B210 -\textgreater{} N210
  (uplink). \texttt{-\/-uplink\ \{wireless,tcp\}} /
  \texttt{-\/-downlink\ \{wireless,tcp\}} lets uplink go over RF while
  the model broadcast returns over plain TCP (the N210 never has to
  transmit). All-TCP mode runs the whole protocol radio-free and was
  validated end-to-end: MNIST 0.92 in 12 rounds.
\item
  \textbf{Waveform choice.} \texttt{-\/-waveform\ sc} uses differential
  DQPSK (default); \texttt{-\/-waveform\ ofdm} uses coherent QPSK, whose
  per-subcarrier pilots track the free-running CFO --- so OFDM-QPSK
  works where single-carrier coherent QPSK would fail.
\end{itemize}

\hypertarget{25-hardware-topology-and-operating-points}{%
\subsubsection{2.5 Hardware topology and operating
points}\label{25-hardware-topology-and-operating-points}}

\begin{itemize}
\tightlist
\item
  \textbf{Radios:} N210 (\texttt{addr=192.168.20.2}, subdev
  \texttt{A:0}, gain 0--31.5 dB, rate 100e6/int) as the access point /
  server; B210s (\texttt{serial=30CD424}, \texttt{30CD3F7}, subdev
  \texttt{A:A}, gain up to \textasciitilde89 dB, flexible rate, 2×2
  MIMO) as the agents / clients.
\item
  \textbf{Rates:} the MARL stack uses a self-consistent 1.6e6 sample /
  0.8e6 symbol default; the N210 needs its exact 2e6/1e6 pairing (1.6e6
  does not snap on the 100 MHz clock).
\item
  \textbf{Operating envelope:} validated per-scheme B210 gains and EVM
  ceilings at 915 MHz; dense constellations (16-QAM and up) need a
  cable, coherent QPSK needs the warm regime.
\end{itemize}

\hypertarget{26-status}{%
\subsubsection{2.6 Status}\label{26-status}}

MARL: per-agent ACK routing, slot sync and logical collision arbitration
are built and sim-validated; hardware pieces (N210 decoded a DQPSK burst
from a B210, \texttt{{[}BURST{]}} line emits and parses) are proven. FL:
compression and the TCP/wireless channel split are done and validated
over TCP; the wireless run is command-ready and waiting on radio time.
This application is mature.

\begin{center}\rule{0.5\linewidth}{0.5pt}\end{center}

\hypertarget{3-application-2--stlc-digital-over-the-air-computation}{%
\subsection{3. Application 2 --- STLC Digital Over-the-Air
Computation}\label{3-application-2--stlc-digital-over-the-air-computation}}

\begin{quote}
Y.-S. Lee, K.-H. Lee, and B. C. Jung, "Space-time coded over-the-air
computation with receive diversity for 6G massive IoT networks,"
\emph{IEEE Wireless Communications Letters}, vol. 15, pp. 31--35, 2026.
(Companion work: "Digital Over-the-Air Computation with Space-Time
Processing for 6G Massive IoT Networks.")
\end{quote}

\hypertarget{31-what-it-is--and-how-it-differs-from-application-1}{%
\subsubsection{3.1 What it is --- and how it differs from Application
1}\label{31-what-it-is--and-how-it-differs-from-application-1}}

\textbf{Over-the-air computation (AirComp)} turns the wireless medium
itself into a calculator: many IoT sensors transmit
\textbf{simultaneously}, their signals \textbf{superimpose} in the air,
and the access point reads a \textbf{function} of them --- the
\textbf{sum / average} --- \emph{without ever decoding an individual
packet}. Where Application 1 asks "did this one device's bytes arrive?",
Application 2 asks "what is the aggregate of all devices' measurements?"
For massive IoT (thousands of sensors reporting to a fusion center) that
is exactly the quantity of interest, and computing it over the air
removes the per-device access bottleneck.

This specific scheme is \textbf{digital} AirComp (values are quantized,
not sent as raw analog amplitudes --- far more noise-robust than
classical analog AirComp) combined with a \textbf{Space-Time Line Code
(STLC)} for \textbf{receive diversity}.

\hypertarget{32-the-scheme}{%
\subsubsection{3.2 The scheme}\label{32-the-scheme}}

\textbf{Each sensor (single transmit antenna, knows only its own
channel):}

\begin{enumerate}
\def\labelenumi{\arabic{enumi}.}
\tightlist
\item
  \textbf{Quantize} its measurement to bits.
\item
  \textbf{Bit-slice} the bits into \textbf{low-order-modulated} symbols
  --- the paper focuses on \textbf{BPSK and DBPSK}.
\item
  \textbf{STLC-encode} across \textbf{2 time slots}, precoding with its
  \emph{own} channel estimate (local CSI at the transmitter).
\item
  \textbf{Transmit} --- time-aligned with every other sensor, so the
  signals \textbf{superimpose} at the access point.
\end{enumerate}

\textbf{Access point (\textgreater= 2 receive antennas):}

\begin{enumerate}
\def\labelenumi{\arabic{enumi}.}
\setcounter{enumi}{4}
\tightlist
\item
  \textbf{Capture} the 2 time slots on the diversity antennas.
\item
  \textbf{Linear-combine} (the STLC combiner) --- a \emph{fixed,
  CSI-free} operation --- to detect the \textbf{superimposed (summed)}
  symbols with full diversity.
\item
  \textbf{De-map / de-bit-slice / de-quantize} to reconstruct the
  \textbf{aggregate}.
\end{enumerate}

\hypertarget{33-why-stlc-and-why-the-receiver-is-nearly-free}{%
\subsubsection{3.3 Why STLC, and why the receiver is nearly
free}\label{33-why-stlc-and-why-the-receiver-is-nearly-free}}

A \textbf{Space-Time Line Code} is the transmit-side dual of the
familiar space-time \emph{block} code (Alamouti). In an STBC the
transmitter has no channel knowledge and the \emph{receiver} does the
channel-aware combining. In an \textbf{STLC the transmitter} uses
\textbf{local} CSI to precode across the 2 time slots so that the
receiver achieves full diversity with a \textbf{simple, CSI-free linear
combine} --- the channel effects cancel in the combining by
construction.

Two consequences make this a good fit for massive IoT and for our
hardware:

\begin{itemize}
\tightlist
\item
  \textbf{Only local CSI at the transmitter} (each sensor its own
  channel) --- no global CSI, low overhead. This is Bang Chul Jung's
  group's signature technique.
\item
  \textbf{The receiver does almost no work.} As the collaborator
  confirmed: \emph{in their frame there is no data processing on the
  receive side except the CSI computation.} The AP estimates CSI (the
  sounding that feeds the transmit precoding), does the fixed linear
  combine, and de-quantizes. There is \textbf{no per-device
  equalization, no FEC, no ARQ} on this path --- the diversity is
  manufactured at the transmitter.
\end{itemize}

\textbf{DBPSK is significant for us:} it is differential, so the AP
needs \textbf{no absolute phase reference}. Given our free-running LOs
(§1.1), this likely removes the shared-10 MHz-reference requirement that
coherent superposition would otherwise impose --- the differential
detection tolerates the independent oscillators. Time alignment across
sensors is still required, and the STLC precoding still needs the local
CSI.

\hypertarget{34-mapping-onto-our-platform}{%
\subsubsection{3.4 Mapping onto our
platform}\label{34-mapping-onto-our-platform}}

Because the receiver is so light and the modulation is low-order, most
of Application 1's heavy machinery (FEC, the full sync/CFO/phase/demod
pipeline, ARQ/CRC) is \textbf{not on this application's critical path}.
What we already have, and what we must add:

\begin{longtable}[]{@{}ll@{}}
\toprule
Need & Status \\
\midrule
\endhead
Low-order modulation (BPSK / DBPSK) & \textbf{have} ---
\texttt{pyphy.modulate} \\
Custom transmit preprocessing (quantize, bit-slice, STLC) &
\textbf{have} --- compose on the \texttt{pyphy} blocks \\
Simultaneous, slot-aligned transmission from N sensors & \textbf{have}
--- reuse \texttt{slot\_sync.py} + \texttt{Radio.transmit} \\
Raw superposition capture at the AP & partial --- \texttt{capture} is
single-antenna today \\
\textbf{Receive diversity (\textgreater= 2 RX antennas)} &
\textbf{build} --- 2-channel capture (\texttt{capture2}) on B210/X310
(both 2×2) \\
Local CSI estimation (channel sounding) & expose --- wrap existing pilot
channel-estimation as a block \\
STLC encode / combine, sum reconstruction & \textbf{build} --- small
numpy blocks \\
\bottomrule
\end{longtable}

So the realistic work is \textbf{one engine addition (a synchronized
2-channel RX capture) + three small numpy blocks (\texttt{stlc\_encode},
\texttt{stlc\_combine}, and the quantize/bit-slice/aggregate app logic)
+ wiring the sounding step}. None of it touches the \texttt{sdr\_system}
decode path; it all lives on the block API of §1.2.

\hypertarget{35-concrete-build-list}{%
\subsubsection{3.5 Concrete build list}\label{35-concrete-build-list}}

\begin{itemize}
\tightlist
\item
  \textbf{Engine:}
  \texttt{Radio.capture2(n)\ -\textgreater{}\ (ant0,\ ant1)} --- a
  synchronized 2-channel UHD RX stream (the one real engine addition;
  receive diversity needs it).
\item
  \textbf{Blocks (numpy):}

  \begin{itemize}
  \tightlist
  \item
    \texttt{stlc\_encode(sym,\ h)} --- device side, 2-slot precoding
    from local CSI.
  \item
    \texttt{stlc\_combine(y0,\ y1{[},\ \ldots{}{]})} --- AP side, the
    fixed CSI-free linear combiner -\textgreater{} superimposed symbol.
  \item
    \texttt{estimate\_channel(rx,\ pilot)\ -\textgreater{}\ h} ---
    expose the existing preamble/pilot estimator.
  \item
    \texttt{quantize} / \texttt{bit\_slice} / \texttt{dequantize} /
    \texttt{aggregate} --- pure-Python application logic.
  \end{itemize}
\item
  \textbf{App loop:} \texttt{slot\_sync} aligns the sensors
  -\textgreater{} each \texttt{stlc\_encode}s its quantized value and
  fires -\textgreater{} AP \texttt{capture2} -\textgreater{}
  \texttt{stlc\_combine} -\textgreater{} de-quantize -\textgreater{} the
  aggregate. Validate a 2-sensor -\textgreater{} 1-AP sum in loopback
  the same way the SC/OFDM chains were validated (0 errors through a
  known offset).
\end{itemize}

\hypertarget{36-open-questions-need-the-papers-equations-to-finalize}{%
\subsubsection{3.6 Open questions (need the paper's equations to
finalize)}\label{36-open-questions-need-the-papers-equations-to-finalize}}

\begin{enumerate}
\def\labelenumi{\arabic{enumi}.}
\tightlist
\item
  \textbf{How is local CSI obtained} --- TDD reciprocity (sensor
  estimates from an AP downlink pilot) or AP-estimates-then-feeds-back?
  This decides whether the AP transmits at all, which matters given our
  RX-only N210 vs the TX-capable B210s.
\item
  \textbf{Number of AP receive antennas M} (diversity order is 2M) ---
  two, or more?
\item
  The exact \textbf{STLC precoding matrix} and \textbf{combiner} (the
  2-slot formulas), the \textbf{quantization + bit-slicing} map, and
  what \textbf{synchronization} the system model assumes.
\end{enumerate}

\hypertarget{37-status}{%
\subsubsection{3.7 Status}\label{37-status}}

Paradigm analyzed and mapped to concrete platform work; the block API is
the right substrate. Not yet implemented --- blocked on the paper's
precise STLC/estimator equations and on building the 2-channel capture.
This is the immediate next application to construct.

\begin{center}\rule{0.5\linewidth}{0.5pt}\end{center}

\hypertarget{4-how-the-two-applications-unify}{%
\subsection{4. How the two applications
unify}\label{4-how-the-two-applications-unify}}

Application 1 fits the existing two archetypes cleanly:

\begin{itemize}
\tightlist
\item
  \textbf{Data-transfer} (FL): \texttt{next\_payload()} +
  \texttt{on\_payload()}.
\item
  \textbf{Control / random-access} (MARL, OSU): \texttt{next\_payload()}
  + \texttt{on\_result(ack)}, ACK = reward.
\end{itemize}

Application 2 does \textbf{not} fit either --- its "output" is neither
delivered bytes nor an ACK, but a \textbf{computed aggregate}, produced
by \textbf{many senders transmitting at once} into a
\textbf{multi-antenna} receiver. It is a \textbf{third archetype:
compute / aggregation.} This is a useful stress test for the UnionLabs
abstraction: the layer must support \emph{simultaneous multi-node
transmission} and \emph{multi-antenna reception}, not only
point-to-point. Concretely, the \texttt{SdrApp} contract extends so that
a sensor contributes an STLC-encoded quantized value and the AP
"receiver" is a combine-and-aggregate step rather than
decode-one-packet.

Stated as one seam:

\begin{longtable}[]{@{}llll@{}}
\toprule
Archetype & Senders & Receiver output & Example \\
\midrule
\endhead
Data-transfer & 1 & the bytes & Federated Learning \\
Control / random-access & 1 of N (contending) & ACK bool (reward) &
MARL, OSU \\
\textbf{Compute / aggregation} & \textbf{N simultaneous} & \textbf{an
aggregate} & \textbf{STLC AirComp} \\
\bottomrule
\end{longtable}

Getting all three behind one interface is the point of the abstraction
layer.

\begin{center}\rule{0.5\linewidth}{0.5pt}\end{center}

\hypertarget{5-reference-appendix}{%
\subsection{5. Reference appendix}\label{5-reference-appendix}}

\hypertarget{51-pyphy-block-reference}{%
\subsubsection{5.1 pyphy block
reference}\label{51-pyphy-block-reference}}

\begin{longtable}[]{@{}ll@{}}
\toprule
Group & Blocks \\
\midrule
\endhead
Framing & \texttt{frame}, \texttt{unframe} \\
FEC & \texttt{fec\_encode}, \texttt{fec\_decode},
\texttt{fec\_decode\_soft}, \texttt{fec\_encoded\_len} (conv \textbar{}
ldpc \textbar{} turbo, k) \\
Modulation & \texttt{modulate}, \texttt{demodulate},
\texttt{soft\_llr} \\
Single-carrier & \texttt{rrc\_tx}, \texttt{rrc\_rx} \\
Sync & \texttt{preamble}, \texttt{acq}, \texttt{cfo\_correct},
\texttt{phase\_correct} \\
OFDM & \texttt{ofdm\_mod}, \texttt{ofdm\_demod},
\texttt{ofdm\_data\_per\_sym} \\
Radio (WITH\_UHD) &
\texttt{Radio(role,args,freq,rate,symbol\_rate,gain,subdev,ant)} ·
\texttt{.transmit()} · \texttt{.capture()} \\
\textbf{To add (App 2)} & \texttt{capture2}, \texttt{stlc\_encode},
\texttt{stlc\_combine}, \texttt{estimate\_channel} \\
\bottomrule
\end{longtable}

Build: \texttt{phy/bindings/build.sh} (DSP only, builds anywhere) or
\texttt{WITH\_UHD=1\ bindings/build.sh} (adds the Radio block, on the
lab host). Worked flowgraph: \texttt{phy/python/phy\_flow\_example.py}.

\hypertarget{52-hardware-inventory}{%
\subsubsection{5.2 Hardware inventory}\label{52-hardware-inventory}}

\begin{longtable}[]{@{}llllll@{}}
\toprule
Device & Role & Address / serial & Subdev & Rate & Notes \\
\midrule
\endhead
N210 & AP / server (RX-only on the rig) & \texttt{addr=192.168.20.2} &
\texttt{A:0} & 2e6/1e6 exact & gain 0--31.5 dB; 100e6/int clock \\
B210 \#1 & agent / client (TX) & \texttt{serial=30CD424} & \texttt{A:A}
& flexible & gain \textasciitilde0--89 dB; 2×2 MIMO \\
B210 \#2 & agent / client (TX) & \texttt{serial=30CD3F7} & \texttt{A:A}
& flexible & 2×2 MIMO (a diversity-RX candidate for App 2) \\
X310 & (optional) & \texttt{addr=} & --- & 200e6/int & 2×2 MIMO \\
\bottomrule
\end{longtable}

\hypertarget{53-where-to-read-the-code-bottom---top}{%
\subsubsection{5.3 Where to read the code (bottom -\textgreater{}
top)}\label{53-where-to-read-the-code-bottom---top}}

\texttt{sdr\_system} (C++ engine, \texttt{-\/-help} surface)
-\textgreater{} \texttt{sdr.py} (auto-generated CLI wrapper)
-\textgreater{} \texttt{marl\_phy.py} / \texttt{real\_channel.py}
(WarmSource / AccessPoint adapters) -\textgreater{}
\texttt{marl\_env.py} (obs / action / reward around the ACK)
-\textgreater{} proposed \texttt{phy\_link.py} (the \texttt{SdrApp} +
\texttt{PayloadSpec} + \texttt{Codec} contract). Full engine reference:
\texttt{../SYSTEM\_REFERENCE.pdf}.

\end{document}
